\chapter{How to Read This Book}
\label{ch:orientation}

This chapter is not about reinforcement learning (RL) itself. It is the map
you will keep coming back to. Read it once carefully now, and skim it again
whenever a later chapter feels like it assumes something you don't
remember — the chances are it's here.

\section{The Mental Model We're Building}

Every chapter in this book is one link in a single chain of reasoning. The
chain looks like this:

\begin{center}
\begin{tikzpicture}[
    node distance=6mm,
    every node/.style={font=\small},
    box/.style={draw, rounded corners, fill=blue!5, align=center,
                minimum width=7.2cm, minimum height=8mm}
]
\node[box] (n1) {Why do we need RL?};
\node[box, below=of n1] (n2) {What makes RL different from supervised learning?};
\node[box, below=of n2] (n3) {Agent, environment, and how they interact};
\node[box, below=of n3] (n4) {Mathematical model: the MDP};
\node[box, below=of n4] (n5) {Measuring how good an action is: return \& value};
\node[box, below=of n5] (n6) {Learning values and policies from experience};
\node[box, below=of n6] (n7) {Scaling up with neural networks (Deep RL)};
\node[box, below=of n7] (n8) {Continuous control for robotics};
\node[box, below=of n8] (n9) {Multiple agents learning together (MARL)};

\foreach \i/\j in {n1/n2, n2/n3, n3/n4, n4/n5, n5/n6, n6/n7, n7/n8, n8/n9}
  \draw[-{Latex[length=2mm]}] (\i) -- (\j);
\end{tikzpicture}
\end{center}

Whenever you feel lost in a later chapter, come back to this diagram and ask:
\emph{which link in this chain am I currently on?} Almost every piece of
confusion in RL comes from mixing up two adjacent links — for example,
confusing ``how good an action is'' (a value) with ``what action to take''
(a policy). We will name and separate these ideas carefully as they appear.

\section{Notation Convention}
\label{sec:notation}

Below is the complete notation table for the whole book. You do not need to
memorize it now — just know it exists, and come back here whenever a symbol
feels unfamiliar. Every symbol below will be re-introduced with a full
explanation the first time it is actually used; this table is a reference,
not a first lesson.

\begin{center}
\begin{tabular}{@{}ll@{}}
\toprule
\textbf{Symbol} & \textbf{Meaning} \\
\midrule
$t$ & discrete time step \\
$\state_t \in \stateSet$ & state at time $t$; the set of all possible states \\
$\obs_t \in \obsSet$ & observation at time $t$ (what the agent actually perceives) \\
$\act_t \in \actSet$ & action at time $t$; the set of all possible actions \\
$\rew_t$ (or $R_t$) & reward received after taking $\act_t$ in $\state_t$ \\
$P(\state' \mid \state, \act)$ & transition dynamics of the environment \\
$\discount \in [0,1]$ & discount factor \\
$\return_t$ & return (discounted sum of future rewards) from time $t$ \\
$\policy(\act \mid \state)$ & stochastic policy; $\policy_\params$ when parameterized \\
$\Vpi(\state)$ & state-value function under policy $\policy$ \\
$\Qpi(\state,\act)$ & action-value function under policy $\policy$ \\
$\Api(\state,\act)$ & advantage function, $\Qpi(\state,\act) - \Vpi(\state)$ \\
$\lr$ & learning rate \\
$\params$ & policy (or, later, network) parameters \\
$\phi$ & critic / value-network parameters (once actor-critic appears) \\
$\replaybuf$ & dataset or replay buffer \\
$\traj$ & trajectory, i.e.\ a sequence $(\state_0, \act_0, \rew_1, \state_1, \dots)$ \\
\bottomrule
\end{tabular}
\end{center}

\begin{remark}
Notation in this book never changes meaning silently. If a later chapter
needs to refine a symbol (for instance, splitting $\params$ into an actor's
$\params$ and a critic's $\phi$), that refinement will be announced
explicitly, in the text, at the moment it happens.
\end{remark}

\section{The Recurring Examples}
\label{sec:recurring-examples}

Instead of inventing a new toy problem every few pages, this book reuses
five environments throughout. Getting familiar with them now means that,
later, you can focus entirely on the \emph{new idea} a chapter introduces,
because the \emph{setting} will already feel like home.

\begin{description}[leftmargin=1.6cm, style=nextline]
  \item[Example A — Grid World] \hfill \\
  A small rectangular grid. The agent starts in one cell and must reach a
  goal cell, choosing among moves like up/down/left/right at each step.
  This is the workhorse example for Chapters 1–12: states, actions,
  rewards, trajectories, policies, value functions, Bellman equations,
  dynamic programming, Monte Carlo, TD learning, and Q-learning are all
  taught on this one small world first.

  \item[Example B — Cliff Walking] \hfill \\
  A grid world variant with a ``cliff'' of very negative reward running
  along one edge. A short, risky path hugs the cliff; a longer, safe path
  avoids it. This tension is exactly what makes it the right environment
  for teaching exploration vs.\ exploitation, and for showing why SARSA
  and Q-learning — two algorithms that look almost identical on paper —
  can learn genuinely different policies (Chapter 11).

  \item[Example C — CartPole] \hfill \\
  A pole balanced on a cart that the agent pushes left or right. The state
  is continuous (position, velocity, angle, angular velocity), which is
  precisely why tables stop working and neural networks enter the story.
  Used from Chapter 13 onward for Deep Q-Networks, policy gradients, and
  PPO.

  \item[Example D — Robotic Arm] \hfill \\
  A simulated arm that must reach, and eventually grasp, a target. Actions
  are continuous torques or joint targets rather than a discrete menu.
  This is where continuous-control algorithms (DDPG, TD3, SAC), reward
  shaping, and sim-to-real transfer are taught, starting in Chapter 22.

  \item[Example E — Multiple Robots] \hfill \\
  Several robots (or arms) that must accomplish a shared task — for
  example, moving objects cooperatively. This is the setting for
  multi-agent RL from Chapter 46 onward: centralized training with
  decentralized execution, credit assignment, and communication.
\end{description}

\begin{intuitionbox}
Think of Examples A and B as your ``training wheels'' world — small enough
to compute everything by hand. Example C is where the book becomes
recognizably modern deep RL. Examples D and E are where the book becomes
\emph{your} field: robotics and multi-robot systems.
\end{intuitionbox}

\section{How Each Concept Will Be Taught}

Every major concept in this book follows the same ten-step shape:
motivation $\to$ intuition $\to$ formal definition $\to$ derivation $\to$
concrete numerical example $\to$ algorithm $\to$ pseudocode $\to$
implementation intuition $\to$ limitations $\to$ connection to what comes
next. You will see this rhythm so often that it becomes a way of reading
new RL material outside of this book too — which is the actual goal.

You will also see four recurring box styles:

\begin{intuitionbox}
Plain-language framing of an idea, before or alongside the mathematics.
\end{intuitionbox}

\begin{keyidea}
A one- or two-sentence distillation of the chapter's central point.
\end{keyidea}

\begin{whythisworks}
An explanation of \emph{why} a method succeeds, not just what it computes.
\end{whythisworks}

\begin{commonmistakes}
Misconceptions beginners commonly form about the concept just introduced.
\end{commonmistakes}

\section{The Code Companion}

At the end of every chapter, after the \emph{Exercises} section, you will
find a short, clearly separated section titled \textbf{Code Companion
Prompts}. These are not part of the textbook's teaching content — they are
ready-to-use prompts (for yourself, or for an AI coding assistant) that
tell you exactly what to implement in the companion repository's
\texttt{code/} directory to turn this chapter's mathematics into running
code. The repository structure is described in the project \texttt{README.md}
and mirrors the recurring examples above: shared environments live in
\texttt{code/envs/}, and each algorithmic chapter gets its own subfolder.

\section{Roadmap Reminder}

Whenever a chapter uses something from earlier without re-deriving it, it
will say so explicitly (``recall from Chapter~5 that\dots''). If you ever
read a sentence in this book that assumes something you don't remember
without pointing back to where it came from, that is a bug in the
book — flag it and we will fix it.

We are now ready to begin the story properly, with the question every RL
course must start from: \emph{why does this field exist at all?}
